%%%%%%%%%%  scan idx 300  |  folio 287  %%%%%%%%%%
Cet homomorphisme est \(K^{\bullet}(Y)\)-linéaire (cf. (2.10)), autrement dit pour \(y\in K^{\bullet}(Y)\) et \(x\in K_{\bullet}(X)\) on a la \emph{formule} dite \emph{de projection}
\begin{equation}\tag{2.11.1.2}
 f_*(yx)=y f_*(x).
\end{equation}
En effet, si \(x=\cl(E)\) pour \(E\in\Ob D^b(X)_{\coh}\) et \(y=\cl(M)\) pour \(M\in\Ob D(Y)_{\uparf}\), on a (III 3.7)
\[
 Rf_*(E)\otimes^{\mathrm L}M\simeq
 Rf_*(E\otimes_Y^{\mathrm L}M).
\]
Enfin, si \(g:Y\to Z\) est un morphisme satisfaisant aux mêmes hypothèses que \(f\) et pour lequel la conjecture de finitude (III 2.1) est vraie, la transitivité des foncteurs images directes implique, pour les homomorphismes (2.11.1.1), la formule
\begin{equation}\tag{2.11.1.3}
 g_*f_*=(gf)_*.
\end{equation}

\textbf{2.11.2.} Soit \(f:X\to Y\) un morphisme propre d'espaces analytiques de dimension finie. Alors, d'après le théorème de finitude de Grauert [2], \(Rf_*\) envoie \(D^b(X)_{\coh}\) dans \(D^b(Y)_{\coh}\), et définit donc un homomorphisme
\begin{equation}\tag{2.11.2.1}
 f_*:K_{\bullet}(X)\longrightarrow K_{\bullet}(Y).
\end{equation}
Comme son analogue algébrique (2.11.1.1), celui-ci est \(K^{\bullet}(Y)\)-linéaire, i.e. la formule (2.11.1.2) est valable, ainsi qu'il résulte de la formule de projection au niveau des catégories dérivées (III 5.9 b)). On a de même la formule de transitivité (2.11.1.3).

\textbf{2.11.3.} Dès qu'on sort du cadre algébrique ou analytique, ce caractère covariant de \(K_{\bullet}\) pour les morphismes propres disparaît. Ainsi, soit \(f:X\to Y\) un morphisme propre d'espaces topologiques annelés par le faisceau des fonctions continues complexes. Alors, si \(E\in\Ob D^b(X)_{\coh}\), on n'a pas en général \(Rf_*(E)\in\Ob D^b(Y)_{\coh}\), comme il est déjà évident pour \(E=\mathcal O_X\) et \(Y\) réduit à un point. En fait, la définition (2.2) de \(K_{\bullet}\)
